\section{对称矩阵与谱定理：什么时候特征方向一定正交}
\label{sec:03-Linear-Algebra/06-Symmetric-and-Positive-Definite-Matrices/01}

上一节我们学会了把矩阵写成 \(A=S\Lambda S^{-1}\)，条件是能找到足够多的特征向量填满 \(S\)。这个分解很强大，但它有两个让人不安的地方。

第一个：\(S\) 的列是特征向量，但它们之间未必正交。两个特征方向可能几乎平行，\(S\) 接近奇异。坐标变换 \(x=Sy\) 会把一个方向上的微小舍入误差放大成另一个方向上的巨大偏差。第二个：实矩阵也可能只有复特征值。旋转矩阵

\[
R=\begin{pmatrix}0&-1\\1&0\end{pmatrix}
\]

的特征值是 \(\pm i\)。你在 \(\R^2\) 里找不到任何一个向量，被 \(R\) 作用后还落在自己张成的直线上。

现在加上一个条件：\(A^{\trans}=A\)。这个条件看起来只是把矩阵沿对角线翻折了一下，没什么戏剧性。但你会发现，它同时解决了上面的两个问题——所有特征值自动变成实数，不同特征值的特征向量自动正交，对角化变成了

\[
A=Q\Lambda Q^{\trans},
\]

\(Q\) 的列是正交归一向量，逆矩阵就是转置。一个不起眼的对称条件，为什么会触发这么强的连锁反应？

\subsection{一条等式锁住了所有方向}

实对称矩阵满足

\[
A^{\trans}=A.
\]

把它塞进一个二次型，立刻得到一个对称关系：

\[
x^{\trans}Ay=(Ax)^{\trans}y.
\]

左边是把 \(y\) 送进 \(A\)、再跟 \(x\) 做内积；右边是把 \(x\) 送进 \(A\)、再跟 \(y\) 做内积。对称性意味着输入与输出的角色可以互换。

这个等式是整个谱定理的支点。下面两个小节分别用它证明"不同特征方向正交"和"特征值不会跑到复平面"。

\subsection{不同特征方向为什么正交}

设 \(Av_i=\lambda_i v_i\)，\(Av_j=\lambda_j v_j\)，并且 \(\lambda_i\ne\lambda_j\)。利用上面的对称关系，把 \(v_i\) 和 \(v_j\) 分别塞进 \(x\) 和 \(y\)：

\[
v_i^{\trans}Av_j=(Av_i)^{\trans}v_j.
\]

左边：\(Av_j=\lambda_j v_j\)，所以是 \(\lambda_j v_i^{\trans}v_j\)。右边：\(Av_i=\lambda_i v_i\)，转置后是 \(\lambda_i v_i^{\trans}v_j\)。两者必须相等：

\[
\lambda_j v_i^{\trans}v_j = \lambda_i v_i^{\trans}v_j.
\]

移项：

\[
(\lambda_j-\lambda_i)v_i^{\trans}v_j=0.
\]

\(\lambda_i\ne\lambda_j\)，所以括号内非零。为此等式成立，只能

\[
v_i^{\trans}v_j=0.
\]

两个属于不同特征值的特征向量，自动正交。没有任何 Gram--Schmidt，没有任何刻意选择——对称性替你做好了这一切。

如果某个特征值重复，对应特征空间可能有多个维度。此时不同特征值的特征向量已经和这个子空间正交，但子空间内部的基还可以任意选取。在这个子空间内部做 Gram--Schmidt，选出一组正交归一基。Gram--Schmidt 的每一步只是已有向量的线性组合，而同一特征值的特征向量的任意线性组合仍是该特征值的特征向量（因为 \(A(\alpha u+\beta v)=\lambda(\alpha u+\beta v)\)），所以正交化不会离开特征空间。最终，整个 \(\R^n\) 可以选出一组完整的正交归一特征基。

\subsection{为什么特征值不会跑到复平面}

即使暂时允许复特征向量，实对称矩阵也会把特征值拉回实数。设 \(Av=\lambda v\)，其中 \(v\) 可能是复向量（\(v\ne 0\)）。用 \(v\) 的共轭转置 \(v^*\) 左乘：

\[
v^*Av = \lambda v^*v.
\]

注意 \(v^*v=\sum|v_i|^2>0\)，所以

\[
\lambda = \frac{v^*Av}{v^*v}.
\]

关键在分子：\(v^*Av\) 是一个标量。它的共轭转置就是它自己——因为 \(A\) 是实对称矩阵，\(A^*=A^{\trans}=A\)，所以

\[
(v^*Av)^* = v^*A^*v = v^*Av.
\]

一个标量等于自己的共轭，它必须是实数。分子是实数，分母是正实数，因此 \(\lambda\) 是实数。不管你怎么尝试塞进复向量，对称矩阵都会把特征值拉回实轴。

\subsection{谱定理}

把上面两条拼在一起，就得到实对称矩阵的谱定理：

\[
A = Q\Lambda Q^{\trans},
\]

其中 \(Q\) 的列是一组正交归一特征向量（\(Q^{\trans}Q=QQ^{\trans}=I\)），\(\Lambda=\diag(\lambda_1,\ldots,\lambda_n)\) 是实对角矩阵。

对比一般对角化 \(A=S\Lambda S^{-1}\)，这里多了三项更强的保证：

\begin{itemize}
\tightlist
\item
  一定存在完整的特征基——对称矩阵永远不会出现 Jordan 块；
\item
  特征基可以选为正交归一的——\(Q^{-1}=Q^{\trans}\)，求逆零成本；
\item
  坐标变换 \(x\mapsto Q^{\trans}x\) 保持长度和夹角——数值上远比 \(S^{-1}\) 稳定。
\end{itemize}

第一条特别值得展开。为什么对称矩阵不会出现 Jordan 块？假设某个特征值 \(\lambda\) 对应了一个大小至少为 2 的 Jordan 块，就存在向量 \(w\) 满足 \((A-\lambda I)w=v\)，其中 \(v\) 是真正的特征向量。但对称性会强制 \(v\) 与 \(w\) 的一个特定组合正交——这跟 Jordan 链的长度假设冲突。详细论证需要用到对称矩阵的 Schur 分解或极小多项式无重复根的性质，但直觉已经清楚：对称性是一个极紧的约束，它不给"半截子"特征方向留生存空间。

\subsection{一个熟悉矩阵的谱分解}

取

\[
A=\begin{pmatrix}2&1\\1&2\end{pmatrix}.
\]

解特征方程：\(\det(A-\lambda I)=(2-\lambda)^2-1=\lambda^2-4\lambda+3=0\)，特征值 \(\lambda_1=3\)，\(\lambda_2=1\)。对应特征向量（未归一）：

\[
\tilde v_1=\begin{pmatrix}1\\1\end{pmatrix},\qquad
\tilde v_2=\begin{pmatrix}1\\-1\end{pmatrix}.
\]

归一化后：

\[
q_1=\frac1{\sqrt2}\begin{pmatrix}1\\1\end{pmatrix},\qquad
q_2=\frac1{\sqrt2}\begin{pmatrix}1\\-1\end{pmatrix}.
\]

于是

\[
Q=\frac1{\sqrt2}
\begin{pmatrix}1&1\\1&-1\end{pmatrix},
\qquad
\Lambda=\begin{pmatrix}3&0\\0&1\end{pmatrix},
\]

验证：\(Q^{\trans}Q=I\)，并且 \(A=Q\Lambda Q^{\trans}\)。

从几何上看这个分解。任意输入向量 \(x\)，先被 \(Q^{\trans}\) 旋转到两条特征轴上——第一条轴沿 \((1,1)\) 方向，第二条轴沿 \((1,-1)\) 方向。沿第一轴拉伸 3 倍，沿第二轴保持不变。最后 \(Q\) 旋转回原坐标。整个过程没有剪切，没有旋转以外的方向混合，只有沿两个正交轴的独立缩放。

\subsection{按方向拆成秩一作用}

把矩阵乘法按列展开：

\[
A = Q\Lambda Q^{\trans}
= \begin{pmatrix}q_1&\cdots&q_n\end{pmatrix}
\begin{pmatrix}\lambda_1&&\\&\ddots&\\&&\lambda_n\end{pmatrix}
\begin{pmatrix}q_1^{\trans}\\\vdots\\q_n^{\trans}\end{pmatrix}
= \sum_{i=1}^{n}\lambda_i\, q_i q_i^{\trans}.
\]

矩阵 \(q_i q_i^{\trans}\) 是投影到第 \(i\) 个特征方向的秩一矩阵。\(A\) 的工作方式可以逐方向阅读：取出 \(x\) 在 \(q_i\) 上的坐标（\(q_i^{\trans}x\)），乘上缩放因子 \(\lambda_i\)，再沿 \(q_i\) 方向输出（乘 \(q_i\)）。然后把所有方向的输出加起来。

这个形式让矩阵函数也一目了然：

\[
f(A)=Q\,f(\Lambda)\,Q^{\trans}
= \sum_{i=1}^{n} f(\lambda_i)\, q_i q_i^{\trans}.
\]

\(f(\Lambda)\) 就是对每个对角元单独作用 \(f\)。平方根、指数、对数、逆——只要 \(f\) 在每个 \(\lambda_i\) 处有定义，就可以直接照搬。\(A^{-1}=Q\Lambda^{-1}Q^{\trans}\)，其中 \(\Lambda^{-1}\) 就是把每个非零特征值取倒数，零特征值对应的维度被丢弃（这就是伪逆的谱版本）。

\subsection{谱定理在数据中的含义}

协方差矩阵 \(\frac1{n}X^{\trans}X\) 是对称半正定矩阵（半正定性在下一节证明）。它的特征向量给出数据变化互不相关的方向——第一个特征向量是方差最大的方向，第二个是在与第一个正交的约束下方差最大的方向，依此类推。特征值给出对应方向上的方差大小。PCA 就是把这些方向按特征值从大到小排列，保留前几个主方向。

同样的结构反复出现。图的 Laplacian 矩阵是对称的，其特征向量给出图上信号的低频与高频模式。核矩阵是对称的，其谱分解给出特征映射的正交坐标。Hessian 矩阵在足够光滑时对称，其正负特征值的个数决定临界点是极小、极大还是鞍点。

对称矩阵出现如此频繁，正是因为自然界中大量"配对关系"天然满足交换对称性：\(x\) 对 \(y\) 的影响方式，与 \(y\) 对 \(x\) 的影响方式，由同一个矩阵刻画。谱定理把这种对偶性翻译成正交方向上的独立缩放。

下一节将沿着谱分解进入一个新的问题：不只看矩阵本身，而是看它产生的标量函数 \(x^{\trans}Ax\)——在什么条件下，这个函数在所有非零方向上恒为正？
